\chapter{监控与告警}
\chapterlead{监控回答两个问题：现在是否正常，异常时如何知道。本章搭一套最小的监控组合：指标、日志、可用性检查与一条能真正送达的告警。}

\section{起步：三条命令判断服务健康}

\begin{shellsession}
$ systemctl --failed
$ df -h / /nix
$ free -h
$ ss -tlnp | grep -c LISTEN
\end{shellsession}

四条命令覆盖四类常见故障：服务崩溃、磁盘满、内存耗尽、端口未监听。把它们写成脚本，定时执行，就是最朴素的监控。

\section{指标：node exporter 与 Prometheus}

\begin{nixcode}
services.prometheus = {
  enable = true;
  exporters.node = {
    enable = true;
    enabledCollectors = [ "systemd" "processes" ];
  };
  scrapeConfigs = [{
    job_name = "node";
    static_configs = [{ targets = [ "127.0.0.1:9100" ]; }];
  }];
};
\end{nixcode}

node exporter 把 CPU、内存、磁盘、网络、systemd 单元状态暴露成文本指标；Prometheus 定时抓取并保存。两者都在本机时配置最简单，对外只开一个端口给 Grafana 或直接查。

\begin{shellsession}
$ curl -s localhost:9100/metrics | grep node_filesystem_avail_bytes | head -3
$ curl -s 'localhost:9090/api/v1/query?query=up'
\end{shellsession}

常用的四条查询值得记住。

\begin{itemize}
  \item \texttt{up}：抓取是否成功。
  \item \texttt{node\_filesystem\_avail\_bytes}：剩余空间。
  \item \texttt{node\_memory\_MemAvailable\_bytes}：可用内存。
  \item \texttt{node\_systemd\_unit\_state}：服务状态。
\end{itemize}

\section{可用性检查}

指标来自服务器自身，服务器整体宕机时它什么都报不出来。外部检查从另一台机器发起。

\begin{nixcode}
services.uptime-kuma.enable = true;
\end{nixcode}

uptime-kuma 提供网页界面，配置 HTTP、TCP、Ping 三类检查，记录可用率与响应时间。它适合放在与业务不同的机器上；只有一台机器时，至少让检查走域名而不是 localhost，这样能覆盖 DNS 与证书问题。

\section{日志}

\begin{nixcode}
services.journald.extraConfig = ''
  SystemMaxUse=1G
  MaxRetentionSec=30day
'';
\end{nixcode}

限制日志占用与保留期，避免日志把磁盘写满。业务日志应当结构化：每行一条记录，含时间、级别、请求标识，便于用 \texttt{journalctl -o json} 与 \texttt{jq} 处理。

\begin{shellsession}
$ journalctl -u api -o json | jq -r 'select(.PRIORITY<"4") | .MESSAGE' | tail -20
\end{shellsession}

这条命令取出该服务最近二十条警告及以上级别的日志，用于快速浏览异常。

\section{一条能送达的告警}

告警的价值取决于能否送达。三条路径，按可靠性排列。

\begin{booktable}{\textwidth}{J{0.769} J{1.231}}
\toprule
方式 & 特点 \\
\midrule
邮件 & 需要 SMTP，容易进垃圾箱 \\
即时消息机器人 & 送达快，需网络可达 \\
短信或电话 & 最可靠，按条计费 \\
\bottomrule
\end{booktable}

个人项目通常用即时消息机器人：systemd 服务里一条 \texttt{curl} 即可。

\begin{nixcode}
systemd.services.alert-disk = {
  script = ''
    avail=$(df --output=avail / | tail -1)
    if [ "$avail" -lt 5242880 ]; then
      curl -s -X POST "$WEBHOOK" -d '{"text":"磁盘剩余不足 5GB"}'
    fi
  '';
  startAt = "*:0/30";
};
\end{nixcode}

告警要有阈值与去重：同一条消息每半小时重复一次就够，不要每次检查都发。长期重复的告警会被忽略，等于没有告警。

\section{练习}

\begin{enumerate}
  \item 写出四条健康检查命令的脚本，输出一行摘要。
  \item 启用 node exporter，用 \texttt{curl} 取出磁盘可用空间的指标并解释数值单位。
  \item 用 Prometheus 查询 \texttt{up}，人为停掉 exporter，观察结果变化。
  \item 配置 journald 的容量与保留期，重建后确认配置生效。
  \item 写一条磁盘告警，阈值设为 5\,GB，验证能收到消息，并说明如何去重。
\end{enumerate}

\begin{checkpoint}
本章结束时，应能：
\begin{itemize}
  \item 用四条命令快速判断服务健康；
  \item 用 node exporter 与 Prometheus 采集与查询指标；
  \item 说明为什么需要外部可用性检查；
  \item 限制日志占用并从结构化的 journal 输出中筛选；
  \item 配置一条能送达、有阈值、有去重的告警。
\end{itemize}
\end{checkpoint}
